\section{Folding preserves the graded codes}

\begin{lemma}[Stability]\label{lem:stab}
Let $1\le i\le m$, $w\in\N^m$, $D\in\N$, $r\in\F$ and $c:=rT^{w_i}$.
\begin{enumerate}
\item[(a)] If $q\in\F[T][x_i,\dots,x_m]_{<k}$ and $\gdeg_w q\le D$, then
$q(c,x_{i+1},\dots,x_m)\in\F[T][x_{i+1},\dots,x_m]_{<k}$ and its graded degree is $\le D$.
\item[(b)] $\Fold_c\big(\Gc_{i-1}^{w,D}\big)\subseteq\Gc_i^{w,D}$.
\item[(c)] For $c_l:=r_lT^{w_l}$ ($1\le l\le i$),
$\Fold_{c_i}\circ\cdots\circ\Fold_{c_1}\big(C_{w,D}\big)$ consists of words of
$\Gc_i^{w,D}$ whose values have $T$-degree at most
$d_i:=\min\big(D,\ (k-1)\sum_{l\le i}w_l\big)$.
\end{enumerate}
\end{lemma}

\begin{proof}
(a) The substitution maps the monomial $T^ax_i^{\,j}x^\beta$ to $r^jT^{a+w_ij}x^\beta$. The
exponents of $x_{i+1},\dots,x_m$ are unchanged, hence $<k$, and the graded degree
$a+w_ij+\sum_{l>i}w_l\beta_l$ is unchanged, hence $\le D$.

(b) Combine (a) with \cref{lem:fold-sub}(b).

(c) By \cref{lem:fold-sub}(b) applied $i$ times the result is $U^{(i)}_q$ with
$q=f(c_1,\dots,c_i,x_{i+1},\dots,x_m)$, which lies in $\Gc_i^{w,D}$ by (a). A monomial
$x^\alpha$ of $f$ contributes the $T$-exponent $\sum_{l\le i}w_l\alpha_l\le(k-1)\sum_{l\le i}w_l$,
and this exponent is at most the graded degree, hence $\le D$.
\end{proof}
\uses{\cref{def:wdeg,def:codes,lem:fold-sub}}\usedby{\cref{prop:complete,lem:round}}
